% Programmation pour la Finance — Version française
% Source de référence: prog_finance.pdf (ne pas modifier le PDF original)
\documentclass[11pt,a4paper]{article}

% Encodage & langue
\usepackage[T1]{fontenc}
\usepackage[utf8]{inputenc}
\usepackage[french]{babel}

% Maths & mise en page
\usepackage{amsmath, amssymb, amsthm}
\usepackage{geometry}
\geometry{margin=2.5cm}
\usepackage{graphicx}
\usepackage{hyperref}
\hypersetup{colorlinks=true, linkcolor=blue, urlcolor=blue, citecolor=blue}
\usepackage{enumitem}

% Titres
\title{Programmation pour la Finance\\\large Version française (synthèse et traduction auteur)}
\author{Cours d'origine: prog\_finance.pdf}
\date{\today}

\begin{document}
\maketitle
\tableofcontents
\newpage

\section*{Notice importante}
\addcontentsline{toc}{section}{Notice importante}
Cette version est une traduction/synthèse en français réalisée à partir du document de référence présent dans le dépôt: \texttt{prog\_finance.pdf}. Le PDF original n'est pas modifié.\newline
\textbf{Objectif:} fournir un support francophone fidèle à l'esprit du cours, sans prétendre reproduire mot-à-mot l'intégralité du texte source.

\paragraph{Méthode.} Nous extrayons la structure et les concepts, puis proposons une reformulation en français. En cas d'écarts d'interprétation, se référer au PDF original.

\section{Introduction}
Cette section présente le contexte du cours et les objectifs d'apprentissage en programmation appliquée à la finance: environnement Python, manipulation de données financières, et mise en œuvre d'algorithmes classiques (p. ex. frontière efficiente).

\subsection{Prérequis}
Compétences de base en Python, calcul matriciel, et notions d'optimisation.

\subsection{Environnement de travail}
Recommandation: Linux, Python 3.10+, gestion d'environnements virtuels, et bibliothèques usuelles (NumPy, Pandas, Matplotlib, CVXOPT/Scipy selon besoins).

\section{Fondamentaux mathématiques}
Rappels synthétiques sur l'algèbre linéaire, statistiques (espérance, variance, covariance), et optimisation convexe utiles à la modélisation de portefeuilles.

\subsection{Notations}
Vecteurs en gras, matrices en capitales, $\mathbb{E}[\cdot]$ pour l'espérance, $\mathrm{Var}(\cdot)$, $\mathrm{Cov}(\cdot,\cdot)$.

\section{Frontière efficiente (aperçu)}
\subsection{Problème de Markowitz}
Formulation quadratique: minimiser la variance $w^\top \Sigma w$ pour un rendement cible $\mu^\top w = r_*$, avec $\sum_i w_i = 1$ et contraintes éventuelles $w_i \ge 0$.

\subsection{Méthodes de résolution}
Approches analytiques (cas sans contraintes), optimisation quadratique numérique, et tracé de la frontière par balayage des rendements cibles.

\subsection{Interprétation}
Compromis risque/rendement, portefeuilles tangents avec un actif sans risque, et sélection du portefeuille selon l'aversion au risque.

\section{Flux de travail recommandé}
\begin{enumerate}[label=\arabic*)]
  \item Préparer un jeu de données (prix/retours) propre et documenté.
  \item Calculer rendements, moyenne vectorielle $\mu$ et matrice de covariance $\Sigma$.
  \item Poser et résoudre le problème d'optimisation.
  \item Tracer la frontière et exporter les résultats.
\end{enumerate}

\section{Références}
\begin{itemize}
  \item PDF source: \texttt{prog\_finance.pdf} (dans le dépôt).
  \item Markowitz, H. (1952). Portfolio Selection.
  \item Outils Python: documentation NumPy, Pandas, SciPy.
\end{itemize}

\end{document}
